\notechapter{N-20220411}{Invertible Matrices, Adjoints, Row Reduction, and Rotations}

\section{Invertibility}

A square matrix \(A\) is invertible if there exists a matrix \(A^{-1}\) such that
\[
  A^{-1}A=I,\qquad AA^{-1}=I.
\]
Some immediate consequences recorded in the note are:
\begin{enumerate}[(i)]
  \item If \(A\) is invertible, then Gaussian elimination has a pivot in every column and every row.
  \item \(A\) may have many distinct powers, but it has only one inverse.
  \item If \(A\) is invertible, then the equation \(Ax=b\) has the unique solution
  \[
    x=A^{-1}b.
  \]
  \item If the homogeneous equation \(Ax=0\) has a nonzero solution, then \(A\) is not invertible.
\end{enumerate}

The last point is fundamental: invertibility is equivalent to trivial kernel for a square matrix.

\section[The 2 x 2 formula]{The \(2\times2\) formula}

For a \(2\times2\) matrix
\[
  A=\begin{pmatrix}a&b\\ c&d\end{pmatrix},
\]
we have
\[
  A^{-1}=\frac1{ad-bc}\begin{pmatrix}d&-b\\ -c&a\end{pmatrix},
\]
provided
\[
  ad-bc\neq0.
\]
Thus the determinant is the obstruction to invertibility.

\section{Diagonal matrices}

If
\[
  A=\operatorname{diag}(d_1,\ldots,d_n),
\]
then \(A\) is invertible iff every \(d_i\neq0\), and
\[
  A^{-1}=\operatorname{diag}(d_1^{-1},\ldots,d_n^{-1}).
\]
This is the simplest case where the inverse is completely visible.

\section{Products of inverses}

For invertible matrices,
\[
  (AB)^{-1}=B^{-1}A^{-1}.
\]
More generally,
\[
  (ABC)^{-1}=C^{-1}B^{-1}A^{-1}.
\]
The order reverses because composition reverses when undoing operations.

\section{Adjugate formula}

Let \(A=(a_{ij})\).  Delete row \(i\) and column \(j\) to get the minor \(M_{ij}\).  The cofactor is
\[
  C_{ij}=(-1)^{i+j}\det M_{ij}.
\]
The adjugate matrix is the transpose of the cofactor matrix:
\[
  \operatorname{adj}(A)=(C_{ji}).
\]
Then
\[
  A^{-1}=\frac1{\det A}\operatorname{adj}(A)
\]
when \(\det A\neq0\).

The note calls this formula computationally expensive but conceptually important.  It proves that entries of the inverse are rational functions of entries of \(A\).

\section{Row reduction for inverses}

The practical method is row reduction:
\[
  [A\mid I]\longrightarrow [I\mid A^{-1}].
\]
For example, for
\[
  M=\begin{pmatrix}1&4\\2&9\end{pmatrix},
\]
one has
\[
  \det M=9-8=1,
\]
and hence
\[
  M^{-1}=\begin{pmatrix}9&-4\\-2&1\end{pmatrix}.
\]
The note also illustrates obtaining this inverse by elementary row operations on the augmented matrix.

\section{Rotation matrices}

The final inserted page discusses planar rotations.  The rotation matrix is
\[
  R_\theta=
  \begin{pmatrix}
  \cos\theta&-\sin\theta\\
  \sin\theta&\cos\theta
  \end{pmatrix}.
\]
It sends
\[
  \begin{pmatrix}x\\y\end{pmatrix}
\]
to
\[
  \begin{pmatrix}
  x\cos\theta-y\sin\theta\\
  x\sin\theta+y\cos\theta
  \end{pmatrix}.
\]
Its inverse is
\[
  R_\theta^{-1}=R_{-\theta}=
  \begin{pmatrix}
  \cos\theta&\sin\theta\\
  -\sin\theta&\cos\theta
  \end{pmatrix}.
\]

\begin{figure}[H]
\centering
\begin{tikzpicture}[scale=1.4,>=Stealth,every node/.style={fill=white,inner sep=1pt}]
  \draw[->] (-0.2,0)--(2.4,0) node[right] {$x$};
  \draw[->] (0,-0.2)--(0,2.05) node[above] {$y$};
  \draw[->,thick] (0,0)--(1.7,0.5);
  \node[right=2pt] at (1.7,0.5) {$v$};
  \draw[->,thick] (0,0)--(1.05,1.42);
  \node[above left=1pt] at (1.05,1.42) {$R_\theta v$};
  \draw (0.55,0.16) arc (16:54:0.57);
  \node at (0.82,0.40) {$\theta$};
\end{tikzpicture}
\caption{A planar rotation acts linearly on vectors and is represented by \(R_\theta\).}
\end{figure}

\section{Expanded synthesis and advanced context}

Invertible matrices are automorphisms of finite-dimensional vector spaces.  Gaussian elimination, determinants, adjugates, and block decompositions are different ways to detect or compute the same object: a reversible linear transformation.


\paragraph{Expanded synthesis.}
For this chapter, the elementary material should be read as a study of what remains invariant when coordinates change. The earlier sections (`Invertibility`, `The \(2\times2\) formula`, `Diagonal matrices`, `Products of inverses`, `Adjugate formula`) compute with arrays, bases, pivots, determinants, or eigenvectors; the deeper object is a linear map together with the spaces it acts on. A matrix is a representation of that map after bases have been chosen. This is why formulas such as
\[
  A' = P^{-1}AP,\qquad \widetilde A = T^{-1}AS
\]
are not cosmetic: they distinguish an intrinsic morphism from a coordinate description.

The structural hierarchy is as follows. Kernels record what a map forgets, images record what it reaches, cokernels record obstructions to solvability, and quotients record information after an equivalence has been imposed. Determinants act on the top exterior power and therefore measure oriented volume; traces are invariant under cyclic rearrangement and become characters in representation theory; adjoints depend on an inner product and lead to self-adjoint spectral theory.

A useful way to return to the computations is to ask, for every formula, which invariant it is measuring. Row reduction measures rank and kernel; diagonalization measures decomposition into invariant subspaces; Gram--Schmidt measures orthogonal projection; positive definiteness measures ellipsoid geometry; tensor notation measures how components transform. The elementary methods are therefore the finite-dimensional laboratory in which duality, quotienting, functoriality, and spectral decomposition can be seen clearly.


\section{Pedagogical repair: inverse as undoing a linear process}

The inverse matrix should be read as an operation that undoes a linear process.  The formula
\[
  (AB)^{-1}=B^{-1}A^{-1}
\]
then becomes obvious: if one first applies \(B\) and then \(A\), one must undo \(A\) first and then undo \(B\).  Row reduction, adjugates, and determinants are three different languages for this same reversibility.

The rotation matrix is a particularly good example because the inverse is visible geometrically.  Rotating by \(\theta\) is undone by rotating by \(-\theta\), so
\[
  R_\theta^{-1}=R_{-\theta}=R_\theta^T.
\]
This is the first glimpse of orthogonal matrices, where algebraic inverse equals geometric transpose.

\section{Elementary-to-advanced bridge: invertible matrices as a Lie group}

The set of all invertible $n\times n$ matrices over $\mathbb R$ is
\[
  GL(n,\mathbb R)=\{A:\det A\ne0\}.
\]
It is a group under multiplication and also an open subset of $\mathbb R^{n^2}$.  Therefore it is a Lie group: a group with a smooth manifold structure.

The tangent space at the identity is the Lie algebra
\[
  \mathfrak{gl}(n,\mathbb R)=M_n(\mathbb R).
\]
The exponential map
\[
  \exp X=I+X+\frac{X^2}{2!}+\cdots
\]
connects infinitesimal generators to invertible transformations.  Elementary matrix powers and exponentials are thus the entry point to continuous symmetry.

\section{Elementary-to-advanced bridge: determinant as a volume form}

The determinant is not merely a formula for invertibility.  It is the factor by which a linear map scales oriented volume:
\[
  \operatorname{Vol}(Av_1,\ldots,Av_n)=\det(A)\operatorname{Vol}(v_1,\ldots,v_n).
\]
The multiplicative identity
\[
  \det(AB)=\det A\det B
\]
says that volume scaling composes multiplicatively.

This is why row operations affect determinants in exactly the familiar ways.  Swapping rows reverses orientation, scaling a row scales volume, and adding one row to another shears volume without changing it.

\section{Elementary-to-advanced bridge: polar decomposition}

Every invertible real matrix has a polar decomposition
\[
  A=QP,
\]
where $Q$ is orthogonal and $P$ is symmetric positive definite.  One can take
\[
  P=(A^TA)^{1/2},\qquad Q=AP^{-1}.
\]
Geometrically, $P$ stretches along orthogonal principal axes, and $Q$ rotates or reflects.  This separates deformation into ``pure strain'' and ``rigid motion.''

The elementary inverse/rotation examples become clearer through this decomposition.  Orthogonal matrices preserve lengths; positive definite matrices perform anisotropic stretching; a general invertible matrix is both.

\section{Elementary-to-advanced bridge: numerical stability of inverses}

Although $A^{-1}$ exists when $\det A\ne0$, computing it may be unstable.  The condition number
\[
  \kappa_2(A)=\|A\|_2\|A^{-1}\|_2=\frac{\sigma_{\max}(A)}{\sigma_{\min}(A)}
\]
measures how much errors may be amplified.  A tiny determinant is not by itself the best diagnostic; singular values describe directional collapse.

Thus the elementary criterion $\det A\ne0$ answers existence, while numerical linear algebra asks sensitivity.

% Second-pass deepening for waves 2-4: wave 2
\section{Second-pass deepening: from row operations to geometry of the general linear group}

Gaussian elimination writes an invertible matrix as a product of elementary matrices.  Geometrically, elementary matrices are simple transformations: shears, scalings, and swaps.  Thus an invertible linear transformation can be built from elementary geometric moves.  This is the computational skeleton behind the group $GL(n)$.

The determinant map
\[
  \det:GL(n,\mathbb R)\to\mathbb R^\times
\]
is a group homomorphism.  Its kernel is the special linear group
\[
  SL(n,\mathbb R)=\{A:\det A=1\}.
\]
This group consists of volume-preserving linear transformations.  The elementary determinant tests in the original note therefore separate invertible maps by how they scale volume.

For rotations, the relevant subgroup is
\[
  O(n)=\{Q:Q^TQ=I\},\qquad SO(n)=O(n)\cap SL(n,\mathbb R).
\]
Orthogonal matrices preserve the inner product, hence all lengths and angles.  The chain
\[
  SO(n)\subset O(n)\subset GL(n)
\]
organizes the original examples: rotations are the most rigid, orthogonal maps allow reflections, and invertible maps allow arbitrary linear distortion without collapse.

\section{Second-pass deepening: exponential map and infinitesimal motion}

A smooth path $A(t)$ in $GL(n)$ with $A(0)=I$ has derivative $X=A'(0)$, an arbitrary matrix.  Conversely, a matrix $X$ generates the path
\[
  A(t)=e^{tX}.
\]
If $X$ is skew-symmetric, then $e^{tX}$ is orthogonal.  Indeed,
\[
  \frac d{dt}(e^{tX})^Te^{tX}=0
\]
when $X^T=-X$.  This gives a precise meaning to ``skew-symmetric matrices generate rotations.''

The elementary rotation matrices in the note are therefore not isolated formulas; they are one-parameter subgroups of $GL(n)$.
